\documentclass[11pt]{article}
\usepackage[a4paper,margin=0.5cm]{geometry} % маленькие поля
\usepackage[utf8]{inputenc} % поддержка UTF-8 
\usepackage[T2A]{fontenc} % кириллические шрифты 
\usepackage[russian]{babel} % правила русского языка
\usepackage{xcolor}   % цвет для выделений
\usepackage{amsmath,amssymb} % формулы

\begin{document}
%\footnotesize % уменьшенный шрифт

 %ВОПРОС: 1
\subsection*{1. Кольцо. Примеры колец. Замкнутость кольца по отношению к операциям
объединения и разности. Алгебра. Утверждение о замкнутости алгебры относительно
дополнения и объединения. Теорема о существовании и единственности
минимального кольца, порожденного системой множеств.} Пусть задано непустое множество $X$, $P(X)$ — семейство всех его подмножеств.\\ \textbf{Определение 2.1.} Непустое семейство $K \subset P(X)$ называют \textbf{кольцом подмножеств}, если для любых $A, B \in K$ выполнены условия: $A \triangle B \in K$, $A \cap B \in K$. \\\textbf{Утверждение 2.1.} Пусть $K \subset P(X)$ — кольцо. Тогда для любых $A, B \in K$ выполняются включения: $A \cup B \in K$, $A \setminus B \in K$. \\\textbf{Доказательство.} Для любых множеств $A$ и $B$ справедливы равенства: $A \cup B = (A \triangle B) \triangle (A \cap B)$ и $A \setminus B = A \triangle (A \cap B)$. Поскольку кольцо замкнуто относительно симметрической разности и пересечения, то $A \cup B \in K$ и $A \setminus B \in K$. \hfill $\Box$ 

 
\textbf{Утверждение 2.2.} Пусть непустая система $K \subset P(X)$ обладает свойствами: 1) для любого $A \in K$ выполнено $X \setminus A \in K$; 2) для любых $A, B \in K$ выполнено $A \cup B \in K$. Тогда $K$ является \textbf{алгеброй}. \\\textbf{Доказательство.} Из условий (1) и (2) получаем: $X = A \cup (X \setminus A) \in K$. Далее: $A \cap B = X \setminus ((X \setminus A) \cup (X \setminus B)) \in K$, $A \setminus B = A \cap (X \setminus B) \in K$, $A \triangle B = (A \cup B) \setminus (A \cap B) \in K$. Таким образом, $K$ — кольцо, содержащее $X$, то есть алгебра. \hfill $\Box$ \\\textbf{Определение 2.3.} Кольцо множеств называется \textbf{$\sigma$-кольцом}, если оно вместе с каждой последовательностью множеств $A_1, A_2, \ldots$ содержит их счётное объединение, то есть $\bigcup_{k=1}^{\infty} A_k \in K$. \\\textbf{Определение 2.4.} \textbf{$\sigma$-алгеброй} называется $\sigma$-кольцо с единицей. 


\textbf{Теорема 2.1.} Для любой непустой системы множеств $S \subset P(X)$ существует единственное минимальное кольцо $K(S)$, содержащее $S$, то есть такое, что: 1) $S \subset K(S)$; 2) для любого кольца $K$, содержащего $S$, выполнено $K(S) \subset K$. \\\textbf{Доказательство.} Существование хотя бы одного кольца, содержащего $S$, очевидно: например, $P(X)$. Рассмотрим множество $M = \bigcup_{A \in S} A$ и кольцо $P(M)$ всех подмножеств $M$. Пусть $\Sigma = \{ K : K \subset P(M), \; S \subset K, \; K \text{ — кольцо} \}$. Положим $K(S) = \bigcap_{K \in \Sigma} K$. Докажем, что $K(S)$ — искомое минимальное кольцо. Во-первых, пересечение любого семейства колец является кольцом (проверка замкнутости относительно $\triangle$ и $\cap$ очевидна). Следовательно, $K(S)$ — кольцо. Во-вторых, $S \subset K(S)$, так как $S$ содержится в каждом кольце из $\Sigma$. Наконец, если $K'$ — произвольное кольцо, содержащее $S$, то $K' \cap P(M)$ является кольцом из $\Sigma$, значит $K(S) \subset K' \cap P(M) \subset K'$. Единственность минимального кольца очевидна. \hfill $\Box$

%ВОРПОС: 2
\subsection*{2. Полукольцо. Примеры полуколец. Лемма о пополнении частичного разбиения до
полного элементами полукольца. Теорема о структуре кольца, минимального над
заданным полукольцом.}
Пусть $X$ — непустое множество, $P(X)$ — семейство всех его подмножеств. \\\textbf{Определение 2.5.} Непустая система $S \subset P(X)$ называется \textbf{полукольцом}, если: 1) $\varnothing \in S$; 2) для любых $A, B \in S$ выполнено $A \cap B \in S$; 3) для любых $A, B \in S$ существует конечная система попарно непересекающихся множеств $C_1, \ldots, C_n \in S$ такая, что $A \setminus B = \bigcup_{k=1}^{n} C_k$. \\\textbf{Лемма 2.1.} Пусть $S$ — полукольцо, $A \in S$, и $A_1, \ldots, A_n \in S$ — попарно непересекающиеся множества, содержащиеся в $A$. Тогда систему $\{A_i\}_{i=1}^{n}$ можно дополнить множествами $A_{n+1}, \ldots, A_m \in S$ до конечного разложения $A = \bigcup_{i=1}^{m} A_i$, причём $A_i \cap A_j = \varnothing$ при $i \ne j$. \\\textbf{Доказательство.} Индукция по $n$. При $n = 1$ утверждение следует из определения полукольца: существуют $C_1, \ldots, C_k \in S$ такие, что $A \setminus A_1 = \bigcup_{j=1}^{k} C_j$, причём множества $C_j$ попарно не пересекаются и не пересекаются с $A_1$. Тогда $A = A_1 \cup C_1 \cup \cdots \cup C_k$ — искомое разложение. Предположим, утверждение верно для $n = k$. Докажем для $n = k + 1$. По предположению индукции существует разложение $A = A_1 \cup \cdots \cup A_k \cup B_1 \cup \cdots \cup B_l$, где все множества в правой части попарно не пересекаются и принадлежат $S$. Рассмотрим множество $A_{k+1}$. Для каждого $j = 1, \ldots, l$ множество $B_j \cap A_{k+1}$ принадлежит $S$ (по замкнутости относительно пересечения). Тогда, по определению полукольца, существуют $D_{j1}, \ldots, D_{jr_j} \in S$, попарно непересекающиеся, такие, что $B_j \setminus A_{k+1} = \bigcup_{s=1}^{r_j} D_{js}$. При этом $B_j = (B_j \cap A_{k+1}) \cup D_{j1} \cup \cdots \cup D_{jr_j}$, и все компоненты не пересекаются. Заметим, что $B_j \cap A_{k+1} \subset A_{k+1}$ и не пересекается с другими $A_i$ при $i \le k$. Теперь заменим в разложении множества $B_j$ на $B_j \cap A_{k+1}, D_{j1}, \ldots, D_{jr_j}$. После проведения такой замены для всех $j$ получим разложение $A$, в котором все множества попарно не пересекаются, принадлежат $S$, и среди них присутствуют $A_1, \ldots, A_{k+1}$. Таким образом, шаг индукции завершён. \hfill $\Box$ 

 
\textbf{Теорема 2.2.} Пусть $S \subset P(X)$ — полукольцо. Тогда минимальное кольцо $K(S)$, порождённое $S$, состоит в точности из всех множеств, представимых в виде конечных объединений попарно непересекающихся элементов из $S$, то есть \[ K(S) = \left\{ A : A = \bigcup_{i=1}^{n} A_i, \; A_i \in S, \; A_i \cap A_j = \varnothing \text{ при } i \ne j \right\}. \] \\\textbf{Доказательство.} Обозначим через $R$ систему всех таких конечных объединений. Покажем, что $R$ — кольцо. 1) Пусть $A, B \in R$. Тогда $A = \bigcup_{i=1}^{n} A_i$, $B = \bigcup_{j=1}^{m} B_j$, где $A_i, B_j \in S$ и внутри каждого объединения множества попарно не пересекаются. Тогда $A \cap B = \bigcup_{1 \le i \le n, 1 \le j \le m} (A_i \cap B_j)$. Каждое $A_i \cap B_j \in S$ (по замкнутости полукольца относительно пересечения), и эти множества попарно не пересекаются (так как не пересекаются $A_i$ между собой и $B_j$ между собой). Следовательно, $A \cap B \in R$. 2) Покажем, что $A \setminus B \in R$. Имеем: \[ A \setminus B = \bigcup_{i=1}^{n} A_i \setminus \left( \bigcup_{j=1}^{m} B_j \right) = \bigcup_{i=1}^{n} \bigcap_{j=1}^{m} (A_i \setminus B_j). \] Для каждого $i$ рассмотрим множества $A_i \setminus B_j$. По определению полукольца, каждое $A_i \setminus B_j$ представляется как конечное объединение непересекающихся элементов из $S$. Затем, применяя лемму, можно показать, что пересечение $\bigcap_{j=1}^{m} (A_i \setminus B_j)$ также допускает представление в виде конечного объединения непересекающихся элементов из $S$. Объединяя такие представления для всех $i$, получим, что $A \setminus B \in R$. 3) Поскольку кольцо замкнуто относительно симметрической разности и пересечения, а также объединения (как следствие), то $R$ действительно является кольцом. Очевидно, что $S \subset R$. Если $K$ — любое кольцо, содержащее $S$, то оно должно содержать все конечные объединения элементов из $S$, то есть $R \subset K$. Следовательно, $R = K(S)$. \hfill $\Box$


%ВОРПОС: 3
\subsection*{3. Мера. Примеры мер. Свойства монотонности и субтрактивности меры.}
\textbf{Определение 3.1.} Пусть на некотором множестве $X$ задано полукольцо $S \subset P(X)$. Будем говорить, что на $S$ задана \textbf{мера}, если каждому элементу $A \in S$ поставлено в соответствие вещественное число $m(A) \in \mathbb{R}$ и при этом выполнены следующие условия: 1) $m(A) \ge 0$ (неотрицательность); 2) если $A = \bigcup_{i=1}^{n} A_i$, где $A_i \in S$ и множества $A_i$ попарно не пересекаются, то $m(A) = \sum_{i=1}^{n} m(A_i)$ (конечная аддитивность). Таким образом, мера первоначально определяется только на полукольце множеств.
\\\textbf{Пример 3.1.} На числовой прямой $\mathbb{R}$ рассмотрим полукольцо $S = \{[a, b) \mid -\infty < a < b < +\infty\}$ и определим меру $m([a, b)) = b - a$. Длина полуинтервала удовлетворяет аксиомам меры. \textbf{Пример 3.2.} На $\mathbb{R}^2$ рассмотрим полукольцо $S$, состоящее из «полуоткрытых» прямоугольников со сторонами, параллельными осям координат: $P = [a_1, b_1) \times [a_2, b_2)$, где $a_i < b_i$. Тогда площадь прямоугольника $m(P) = (b_1 - a_1)(b_2 - a_2)$ удовлетворяет аксиомам меры.
 \\\textbf{Определение 3.2.} Мера $m$, заданная на полукольце $S \subset P(X)$, называется \textbf{счётно-аддитивной} (или \textbf{$\sigma$-аддитивной}), если для любой последовательности попарно непересекающихся множеств $A_1, A_2, \ldots \in S$ такой, что $\bigcup_{i=1}^{\infty} A_i \in S$, выполнено равенство $m\left( \bigcup_{i=1}^{\infty} A_i \right) = \sum_{i=1}^{\infty} m(A_i)$.\\ \textbf{Свойство 3.1 (Монотонность).} Если $A, B \in K$ и $A \subset B$, то $m(A) \le m(B)$. \\\textbf{Доказательство.} Поскольку $B = A \cup (B \setminus A)$ и $A \cap (B \setminus A) = \varnothing$, причём $B \setminus A \in K$ (кольцо замкнуто относительно разности), то по аддитивности меры $m(B) = m(A) + m(B \setminus A) \ge m(A)$, так как $m(B \setminus A) \ge 0$. \hfill $\Box$ \\\textbf{Свойство 3.2 (Субтрактивность меры).} Если $A, B \in K$, $A \subset B$ и $m(A) < \infty$, то $m(B \setminus A) = m(B) - m(A)$. \\\textbf{Доказательство.} Из равенства $m(B) = m(A) + m(B \setminus A)$ следует искомая формула, поскольку все величины конечны. \hfill $\Box$

%ВОРПОС: 4
\subsection*{4. Свойства меры: мера объединения, симметрической разности, оценка модуля
разности мер множеств, оценка меры симметрической разности двух множеств,
счетная полуаддитивность меры.}
\textbf{Свойство 3.4 (Включение–исключение для двух множеств).} Если $A, B \in K$, то $m(A \cup B) = m(A) + m(B) - m(A \cap B)$. \\\textbf{Доказательство.} Запишем $A \cup B$ как объединение непересекающихся множеств: $A \cup B = (A \setminus B) \cup (A \cap B) \cup (B \setminus A)$. Тогда по аддитивности $m(A \cup B) = m(A \setminus B) + m(A \cap B) + m(B \setminus A)$. С другой стороны, $A = (A \setminus B) \cup (A \cap B)$, откуда $m(A) = m(A \setminus B) + m(A \cap B)$. Аналогично, $m(B) = m(B \setminus A) + m(A \cap B)$. Выражая $m(A \setminus B)$ и $m(B \setminus A)$ и подставляя в первое равенство, получаем требуемую формулу. \hfill $\Box$ \\\textbf{Свойство 3.5 (Мера симметрической разности).} Если $A, B \in K$, то $m(A \triangle B) = m(A) + m(B) - 2m(A \cap B)$. \textbf{Доказательство.} Так как $A \triangle B = (A \setminus B) \cup (B \setminus A)$ и множества не пересекаются, то $m(A \triangle B) = m(A \setminus B) + m(B \setminus A)$. Используя равенства $m(A \setminus B) = m(A) - m(A \cap B)$ и $m(B \setminus A) = m(B) - m(A \cap B)$ (которые справедливы, если мера конечна, либо в общем случае с осторожностью), получаем искомую формулу. \hfill $\Box$ \\\textbf{Свойство 3.6 (Неравенство для разности мер).} Если $A, B \in K$, то $|m(A) - m(B)| \le m(A \triangle B)$. \textbf{Доказательство.} Заметим, что $A \subset B \cup (A \triangle B)$, откуда по монотонности и полуаддитивности $m(A) \le m(B) + m(A \triangle B)$. Аналогично, $B \subset A \cup (A \triangle B)$, поэтому $m(B) \le m(A) + m(A \triangle B)$. Из этих двух неравенств следует требуемое. \hfill $\Box$ \\\textbf{Свойство 3.7 (Неравенство треугольника для симметрической разности).} Если $A, B, C \in K$, то $m(A \triangle B) \le m(A \triangle C) + m(C \triangle B)$. \textbf{Доказательство.} Используем включение $A \triangle B \subset (A \triangle C) \cup (C \triangle B)$. Действительно, если $x \in A \triangle B$, то $x$ принадлежит ровно одному из множеств $A$ и $B$. Если $x \in C$, то $x \notin B$ (иначе $x$ принадлежал бы обоим $A$ и $B$), значит, $x \in C \triangle B$. Если же $x \notin C$, то $x \in A \triangle C$. Таким образом, включение доказано. Тогда по монотонности и полуаддитивности $m(A \triangle B) \le m((A \triangle C) \cup (C \triangle B)) \le m(A \triangle C) + m(C \triangle B)$. \hfill $\Box$ \\\textbf{Свойство 3.8 (Счётная полуаддитивность).} Пусть мера $m$ является $\sigma$-аддитивной на кольце $K$. Если $A_1, A_2, \ldots \in K$ (могут перескаться!!!) и $\bigcup_{i=1}^{\infty} A_i \in K$, то $m\!\left(\bigcup_{i=1}^{\infty} A_i\right) \le \sum_{i=1}^{\infty} m(A_i)$. \textbf{Доказательство.} Построим последовательность попарно непересекающихся множеств $B_k \in K$ следующим образом: $B_1 = A_1$, $B_k = A_k \setminus \bigcup_{i=1}^{k-1} A_i$ для $k \ge 2$. Так как кольцо замкнуто относительно конечных объединений и разностей, то $B_k \in K$. Кроме того, $B_k \subset A_k$ и $\bigcup_{k=1}^{\infty} B_k = \bigcup_{k=1}^{\infty} A_k = A$. В силу $\sigma$-аддитивности меры $m(A) = \sum_{k=1}^{\infty} m(B_k) \le \sum_{k=1}^{\infty} m(A_k)$, где неравенство следует из монотонности ($m(B_k) \le m(A_k)$). \hfill $\Box$

%ВОРПОС: 5
\subsection*{5. Теорема о счетной аддитивности полуинтервальной меры. Пример не sigma-аддитивной
меры на числовой прямой. Примеры sigma-аддитивных мер.}
\textbf{Теорема 3.1 (3.1).} Длина полуинтервала, определённая на полукольце $S = \{[a, b) \subset \mathbb{R}\}$ формулой $m([a, b)) = b - a$, является $\sigma$-аддитивной мерой. \\\textbf{Доказательство.} Пусть $A = [a, b) \in S$ и $A = \bigcup_{i=1}^{\infty} A_i$, где $A_i = [a_i, b_i) \in S$ и множества $A_i$ попарно не пересекаются. Требуется доказать, что $b - a = \sum_{i=1}^{\infty} (b_i - a_i)$. 1. Доказательство неравенства $\sum_{i=1}^{\infty}(b_i - a_i) \le b - a$. Для любого натурального $n$ имеем $\bigcup_{i=1}^{n} A_i \subset A$, и так как мера аддитивна для конечных объединений, то $\sum_{i=1}^{n} (b_i - a_i) = \sum_{i=1}^{n} m(A_i) = m(\bigcup_{i=1}^{n} A_i) \le m(A) = b - a$. Переходя к пределу при $n \to \infty$, получаем $\sum_{i=1}^{\infty}(b_i - a_i) \le b - a$. Ряд с неотрицательными членами сходится. 2. Доказательство неравенства $b - a \le \sum_{i=1}^{\infty}(b_i - a_i)$. Зафиксируем произвольное $\varepsilon > 0$. Для каждого полуинтервала $A_i = [a_i, b_i)$ построим интервал $B_i = (a_i - \tfrac{\varepsilon}{2^{i+1}}, b_i)$, так что $A_i \subset B_i$. Также вместо полуинтервала $A$ рассмотрим отрезок $B = [a, b - \tfrac{\varepsilon}{2}] \subset A$. Тогда $B \subset A = \bigcup_{i=1}^{\infty} A_i \subset \bigcup_{i=1}^{\infty} B_i$. Таким образом, отрезок $B$ покрыт системой открытых интервалов $\{B_i\}$. По лемме Гейне–Бореля из этого покрытия можно выделить конечное подпокрытие: существуют индексы $i_1, \ldots, i_k$ такие, что $B \subset \bigcup_{j=1}^{k} B_{i_j}$. Тогда длина отрезка $B$ не превосходит суммы длин интервалов этого конечного покрытия: $b - a - \tfrac{\varepsilon}{2} = |B| \le \sum_{j=1}^{k} |B_{i_j}| = \sum_{j=1}^{k} (b_{i_j} - a_{i_j} + \tfrac{\varepsilon}{2^{i_j+1}}) \le \sum_{i=1}^{\infty} (b_i - a_i) + \sum_{i=1}^{\infty} \tfrac{\varepsilon}{2^{i+1}}$. Поскольку $\sum_{i=1}^{\infty} \tfrac{\varepsilon}{2^{i+1}} = \tfrac{\varepsilon}{2}$, получаем $b - a - \tfrac{\varepsilon}{2} \le \sum_{i=1}^{\infty} (b_i - a_i) + \tfrac{\varepsilon}{2}$, откуда $b - a \le \sum_{i=1}^{\infty} (b_i - a_i) + \varepsilon$. В силу произвольности $\varepsilon$ заключаем, что $b - a \le \sum_{i=1}^{\infty}(b_i - a_i)$. Из двух доказанных неравенств следует равенство (3.1), что и доказывает $\sigma$-аддитивность длины. \hfill $\Box$

\textbf{Пример 3.5.} Пусть $x_0$ — фиксированная точка множества $X$. Для любого множества $A \subset X$ положим $\delta_{x_0}(A) = 1$, если $x_0 \in A$, и $\delta_{x_0}(A) = 0$, если $x_0 \notin A$. Эта мера, называемая \textbf{мерой Дирака} (или единичной массой в точке $x_0$), является $\sigma$-аддитивной на $S = P(X)$.

%ВОРПОС: 6
\subsection*{6. Продолжение меры. Теорема о продолжении меры, заданной на полукольце, на
минимальное кольцо. Следствие о сохранении свойства продолжения меры}
\textbf{Определение 4.1.} Пусть $m$ — мера, заданная на полукольце $S$, и пусть $K$ — кольцо, содержащее $S$ (т.е. $S \subset K$). Мера $\mu$, заданная на $K$, называется \textbf{продолжением меры} $m$, если для любого множества $A \in S$ выполняется равенство $\mu(A) = m(A)$. \\\textbf{Теорема 4.1.} Пусть $m$ — мера на полукольце $S \subset P(X)$ и $K(S)$ — минимальное кольцо, порождённое $S$. Тогда на $K(S)$ существует единственная мера $\mu$, являющаяся продолжением меры $m$. \\\textbf{Доказательство.} Доказательство разобьём на три этапа: построение продолжения, проверка корректности построения и проверка аксиом меры. 1. \textit{Построение.Доказательство ЕДИНСТВЕННОСТИ (если продолжение существует)} В силу теоремы о структуре минимального кольца, порождённого полукольцом (теорема 2.2), каждое множество $A \in K(S)$ допускает представление в виде конечного объединения попарно непересекающихся элементов из $S$: $A = \bigcup_{i=1}^{n} A_i$, где $A_i \in S$, $A_i \cap A_j = \varnothing$ при $i \ne j$. Если на $K(S)$ существует продолжение $\mu$ меры $m$, то в силу аддитивности должно выполняться: $\mu(A) = \sum_{i=1}^{n} \mu(A_i) = \sum_{i=1}^{n} m(A_i) (1)$. Это равенство задаёт явную формулу для продолжения, если оно существует. Таким образом, продолжение, если оно существует, единственно. 2. \textit{Корректность.} Покажем, что формула (1) задаёт функцию $\mu : K(S) \to [0, +\infty]$ корректно, т.е. значение $\mu(A)$ не зависит от выбора разбиения множества $A$ на попарно непересекающиеся элементы из $S$. Пусть $A \in K(S)$ имеет два разложения: $A = \bigcup_{i=1}^{n} A_i = \bigcup_{j=1}^{m} B_j$, где $A_i, B_j \in S$ и внутри каждого разложения множества попарно не пересекаются. Рассмотрим множества $C_{ij} = A_i \cap B_j$. Так как $S$ — полукольцо, то $C_{ij} \in S$. Для каждого фиксированного $i$ множества $\{C_{ij}\}_{j=1}^{m}$ попарно не пересекаются и $A_i = \bigcup_{j=1}^{m} C_{ij}$. Аналогично, для каждого фиксированного $j$ множества $\{C_{ij}\}_{i=1}^{n}$ попарно не пересекаются и $B_j = \bigcup_{i=1}^{n} C_{ij}$. Тогда, используя конечную аддитивность меры $m$ на $S$, получаем: $\sum_{i=1}^{n} m(A_i) = \sum_{i=1}^{n} \sum_{j=1}^{m} m(C_{ij}) = \sum_{j=1}^{m} \sum_{i=1}^{n} m(C_{ij}) = \sum_{j=1}^{m} m(B_j)$. Таким образом, значение $\mu(A)$ не зависит от выбора разложения, и функция $\mu$ определена корректно. 3. \textit{Проверка аксиом меры.} • Неотрицательность: по определению, $\mu(A) = \sum_{i=1}^{n} m(A_i) \ge 0$, так как $m(A_i) \ge 0$. • Аддитивность: пусть $A, B \in K(S)$ и $A \cap B = \varnothing$. Тогда существуют разложения $A = \bigcup_{i=1}^{n} A_i$, $B = \bigcup_{j=1}^{m} B_j$, где $A_i, B_j \in S$ и множества внутри каждого объединения попарно не пересекаются. Поскольку $A$ и $B$ не пересекаются, объединение $A \cup B$ можно представить как $\bigcup_{i=1}^{n} A_i \cup \bigcup_{j=1}^{m} B_j$, где все множества попарно не пересекаются. Тогда по определению $\mu$: $\mu(A \cup B) = \sum_{i=1}^{n} m(A_i) + \sum_{j=1}^{m} m(B_j) = \mu(A) + \mu(B)$. Это доказывает конечную аддитивность $\mu$. Более того, если $A_1, \ldots, A_k \in K(S)$ попарно не пересекаются, то аналогичным образом можно показать, что $\mu(\bigcup_{p=1}^{k} A_p) = \sum_{p=1}^{k} \mu(A_p)$. Таким образом, $\mu$ — мера на $K(S)$, которая по построению совпадает с $m$ на $S$. Единственность уже обоснована на этапе построения. Теорема доказана. \hfill $\Box$

%ВОРПОС: 7
\subsection*{7. Множество меры нуль. Примеры множеств меры нуль. Непрерывность сверху (снизу)
меры на полукольце. Теорема о связи счетной аддитивности и непрерывности меры.}
\textbf{Определение 4.2.} Пусть на кольце $K \subset P(X)$ задана $\sigma$-аддитивная мера $\mu$. Множество $A \subset X$ (не обязательно принадлежащее $K$) называется множеством меры нуль, если для любого $\varepsilon > 0$ существует конечный или счётный набор множеств $E_i \in K$, покрывающий $A$ (т.е. $A \subset \bigcup_{i=1}^\infty E_i$), такой что $\sum_{i=1}^\infty \mu(E_i) < \varepsilon$. \\ \textbf{Определение 5.1 (Непрерывность снизу).} Меру $m$, заданную на кольце $K$, называют непрерывной снизу, если для любой возрастающей последовательности множеств $A_1 \subseteq A_2 \subseteq \dots$, такой что $A_n \in K$ и $A = \bigcup_{n=1}^\infty A_n \in K$, выполняется равенство $\lim_{n \to \infty} m(A_n) = m(A)$. Обозначение: $A_n \uparrow A$. \\ \textbf{Определение 5.2 (Непрерывность сверху).} Меру $m$, заданную на кольце $K$, называют непрерывной сверху, если для любой убывающей последовательности множеств $A_1 \supseteq A_2 \supseteq \dots$, такой что $A_n \in K$, $A = \bigcap_{n=1}^\infty A_n \in K$ и $m(A_1) < \infty$, выполняется равенство $\lim_{n \to \infty} m(A_n) = m(A)$. Обозначение: $A_n \downarrow A$. Условие $m(A_1) < \infty$ гарантирует, что все меры конечны (в силу монотонности) и разности $m(A_n) - m(A)$ имеют смысл. \\ \textbf{Теорема 5.1.} Мера $m$, заданная на кольце $K$, является счётно-аддитивной тогда и только тогда, когда она непрерывна снизу. Кроме того, если мера счётно-аддитивна и $A_n \downarrow A$ с $m(A_1) < \infty$, то она непрерывна сверху. \\ \textbf{Доказательство.} Доказательство разбиваем на две части: необходимость и достаточность. \textbf{Необходимость.} Предположим, что мера $m$ счётно-аддитивна на кольце $K$. Докажем непрерывность снизу. Пусть $A_n \uparrow A$, где $A_n, A \in K$. Построим последовательность попарно непересекающихся множеств: $B_1 = A_1$, $B_n = A_n \setminus A_{n-1}$ для $n \geq 2$. Так как $K$ — кольцо, то $B_n \in K$. Кроме того, $A = \bigcup_{n=1}^\infty A_n = \bigsqcup_{n=1}^\infty B_n$. В силу счётной аддитивности имеем $m(A) = \sum_{n=1}^\infty m(B_n)$. С другой стороны, для любого $n \geq 1$ выполнено $A_n = \bigcup_{k=1}^n B_k$, поэтому $m(A_n) = \sum_{k=1}^n m(B_k)$. Таким образом, частичные суммы ряда $\sum_{k=1}^\infty m(B_k)$ равны $m(A_n)$, и сходимость ряда влечёт $\lim_{n \to \infty} m(A_n) = \sum_{k=1}^\infty m(B_k) = m(A)$. Значит, мера непрерывна снизу. Теперь докажем непрерывность сверху при дополнительном условии конечности меры первого множества. Пусть $A_n \downarrow A$, $A_n, A \in K$ и $m(A_1) < \infty$. Рассмотрим возрастающую последовательность множеств $B_n = A_1 \setminus A_n$. Так как $A_n$ убывает, то $B_n$ возрастает. Кроме того, $\bigcup_{n=1}^\infty B_n = A_1 \setminus A$. По только что доказанному свойству непрерывности снизу для последовательности $B_n \uparrow (A_1 \setminus A)$ имеем $\lim_{n \to \infty} m(B_n) = m(A_1 \setminus A)$. Но $m(B_n) = m(A_1) - m(A_n)$ и $m(A_1 \setminus A) = m(A_1) - m(A)$. Следовательно, $\lim_{n \to \infty} (m(A_1) - m(A_n)) = m(A_1) - m(A)$, откуда $\lim_{n \to \infty} m(A_n) = m(A)$. Таким образом, мера непрерывна сверху. \textbf{Достаточность}. Предположим теперь, что мера $m$ непрерывна снизу. Докажем её счётную аддитивность. Пусть $\{B_i\}_{i=1}^\infty$ — последовательность попарно непересекающихся множеств из $K$ таких, что $B = \bigsqcup_{i=1}^\infty B_i \in K$. Определим возрастающую последовательность множеств $A_n = \bigcup_{i=1}^n B_i \in K$. Очевидно, $A_n \uparrow B$. По условию непрерывности снизу имеем $\lim_{n \to \infty} m(A_n) = m(B)$. Но $m(A_n) = \sum_{i=1}^n m(B_i)$ в силу конечной аддитивности меры. Поэтому $m(B) = \lim_{n \to \infty} \sum_{i=1}^n m(B_i) = \sum_{i=1}^\infty m(B_i)$, что и означает счётную аддитивность меры $m$. \hfill $\Box$

%ВОРПОС: 8
\subsection*{8. Внешняя мера. Утверждение о внешней мере. Свойства внешней меры: мера
элементов кольца, неотрицательность, монотонность.}
Пусть на алгебре $K$ задана $\sigma$-аддитивная мера $m$. Определим на системе всех подмножеств множества $X$ функцию $\mu^{*}$, которая каждому множеству $A \subset X$ ставит в соответствие число $\mu^{*}(A)$ по следующему правилу. Так как $X \in K$, то для любого множества $A$ существует его покрытие элементами алгебры $K$ (например, $A_{1} = X$, $A_{i} = \varnothing$ при $i \geq 2$). Вычислим меру такого покрытия и возьмём точную нижнюю грань по всевозможным покрытиям. 
\\\textbf{Определение}Внешней мерой множества $A \subset X$ называется число $\mu^{*}(A) = \inf \{ \sum_{i=1}^{\infty} m(A_{i}) : A \subset \bigcup_{i=1}^{\infty} A_{i}, \; A_{i} \in K \}$. \\\textbf{Теорема 6.1.} Внешняя мера $\mu^{*}$, заданная на $P(X)$, является продолжением меры $m$ с алгебры $K$, то есть для любого множества $A \in K$ выполняется равенство $\mu^{*}(A) = m(A)$. \textbf{Доказательство.} Пусть $A \in K$. С одной стороны, рассмотрим покрытие $A_{1} = A$, $A_{i} = \varnothing$ при $i \geq 2$. Тогда $\mu^{*}(A) \leq \sum_{i=1}^{\infty} m(A_{i}) = m(A)$. С другой стороны, пусть $\{A_{i}\}_{i=1}^{\infty}$ — произвольное покрытие множества $A$ элементами из $K$, то есть $A \subset \bigcup_{i=1}^{\infty} A_{i}$. Поскольку мера $m$ $\sigma$-аддитивна на алгебре $K$, она обладает свойством счётной полуаддитивности: $m(A) \leq \sum_{i=1}^{\infty} m(A_{i})$. Таким образом, $m(A)$ является нижней границей для сумм $\sum_{i=1}^{\infty} m(A_{i})$ по всем покрытиям. Следовательно, $m(A) \leq \inf \{ \sum_{i=1}^{\infty} m(A_{i}) : A \subset \bigcup_{i=1}^{\infty} A_{i}, \; A_{i} \in K \} = \mu^{*}(A)$. Из двух неравенств получаем $\mu^{*}(A) = m(A)$. \\\textbf{Свойство 6.1 (Совпадение с мерой на алгебре).} Если $A \in K$, то $\mu^{*}(A) = m(A)$. \\\textbf{Свойство 6.2 (Неотрицательность).} Для любого множества $A \subset X$ выполнено $\mu^{*}(A) \geq 0$, а также $\mu^{*}(\varnothing) = 0$. Действительно, по определению внешняя мера есть инфимум неотрицательных чисел, поэтому $\mu^{*}(A) \geq 0$.\\\textbf{Свойство 6.3 (Монотонность).} Если $A, B \subset X$ и $A \subseteq B$, то $\mu^{*}(A) \leq \mu^{*}(B)$. Действительно, любое покрытие множества $B$ элементами из $K$ является также покрытием для $A$. Поэтому множество сумм, по которому берётся инфимум для $\mu^{*}(A)$, содержит множество сумм для $\mu^{*}(B)$. Следовательно, инфимум для $A$ не превосходит инфимума для $B$, то есть $\mu^{*}(A) \leq \mu^{*}(B)$.

%ВОРПОС: 9
\subsection*{9. Свойства внешней меры: счетная полуаддитивность, оценка модуля разности мер,
оценка меры симметрической разности двух множеств. Внутренняя мера.
Соотношение внутренней и внешней меры одного и того же множества.}
\textbf{Свойство 6.4 (Счётная полуаддитивность).} Для любой последовательности множеств $\{B_i\}_{i=1}^\infty \subset X$ выполняется неравенство $\mu^{*}(\bigcup_{i=1}^\infty B_i) \leq \sum_{i=1}^\infty \mu^{*}(B_i)$. Доказательство. Если ряд справа расходится, то неравенство очевидно. Предположим, что ряд сходится. Зафиксируем произвольное $\varepsilon > 0$. По определению внешней меры для каждого $i$ существует покрытие $\{A_{ij}\}_{j=1}^\infty$ множества $B_i$ элементами из $K$ такое, что $\sum_{j=1}^\infty m(A_{ij}) \leq \mu^{*}(B_i) + \frac{\varepsilon}{2^i}$. Объединение всех этих покрытий образует покрытие множества $\bigcup_{i=1}^\infty B_i$. Тогда по определению внешней меры имеем $\mu^{*}(\bigcup_{i=1}^\infty B_i) \leq \sum_{i=1}^\infty \sum_{j=1}^\infty m(A_{ij}) \leq \sum_{i=1}^\infty (\mu^{*}(B_i) + \frac{\varepsilon}{2^i}) = \sum_{i=1}^\infty \mu^{*}(B_i) + \varepsilon$. Так как $\varepsilon$ произвольно, получаем требуемое неравенство. \\\textbf{Свойство 6.5 (Неравенство для разности внешних мер).} Для любых множеств $A, B \subset X$ выполняется $|\mu^{*}(A) - \mu^{*}(B)| \leq \mu^{*}(A \triangle B)$. Доказательство. Заметим, что $A \subset B \cup (A \triangle B)$. По свойствам монотонности и счётной полуаддитивности имеем $\mu^{*}(A) \leq \mu^{*}(B) + \mu^{*}(A \triangle B)$. Аналогично, $B \subset A \cup (A \triangle B)$, откуда $\mu^{*}(B) \leq \mu^{*}(A) + \mu^{*}(A \triangle B)$. Из этих двух неравенств следует требуемое. \\\textbf{Свойство 6.6 (Неравенство треугольника для симметрической разности).} Для любых множеств $A, B, C \subset X$ выполняется $\mu^{*}(A \triangle B) \leq \mu^{*}(A \triangle C) + \mu^{*}(C \triangle B)$. Доказательство. Используем включение $A \triangle B \subset (A \triangle C) \cup (C \triangle B)$. Действительно, если $x \in A \triangle B$, то $x$ принадлежит ровно одному из множеств $A$ и $B$. Если $x \in C$, то $x \notin B$, значит $x \in C \triangle B$. Если же $x \notin C$, то $x \in A \triangle C$. Таким образом, включение доказано. Тогда по монотонности и счётной полуаддитивности имеем $\mu^{*}(A \triangle B) \leq \mu^{*}((A \triangle C) \cup (C \triangle B)) \leq \mu^{*}(A \triangle C) + \mu^{*}(C \triangle B)$. \\\textbf{Определение 6.2.} Внутренней мерой множества $A \subset X$ называется число $\mu_{*}(A) = m(X) - \mu^{*}(X \setminus A)$. \\\textbf{Свойство 6.7 (Связь внутренней и внешней мер).} Для любого множества $A \subset X$ выполняется неравенство $\mu_{*}(A) \leq \mu^{*}(A)$. Доказательство. Запишем очевидное равенство $X = A \cup (X \setminus A)$. Применяя свойство счётной полуаддитивности внешней меры, получаем $m(X) = \mu^{*}(X) \leq \mu^{*}(A) + \mu^{*}(X \setminus A)$. Отсюда $m(X) - \mu^{*}(X \setminus A) \leq \mu^{*}(A)$. Но левая часть равна $\mu_{*}(A)$ по определению. Таким образом, $\mu_{*}(A) \leq \mu^{*}(A)$.

%ВОРПОС: 10
\subsection*{10. Множество, измеримое по Лебегу относительно меры. Измеримость по Лебегу
множеств меры нуль. Критерий измеримости. Следствие о мере симметрической
разности с измеримым множеством.}
Пусть $X$ — произвольное множество, $K \subset P(X)$ — алгебра его подмножеств, на которой задана $\sigma$-аддитивная конечная мера $m$ (т.е. $m(X) < \infty$). Для каждого множества $A \subset X$ определена внешняя мера $\mu^{*}$. \\\textbf{Определение 7.1.} Множество $A \subset X$ называется измеримым по Лебегу относительно меры $m$, если для него выполнено равенство $\mu^{*}(A) + \mu^{*}(X \setminus A) = m(X)$. \\\textbf{Упражнение 7.1.} Доказать, что множество $A$ измеримо тогда и только тогда, когда $\mu^{*}(A) = \mu_{*}(A)$, где $\mu_{*}(A) = m(X) - \mu^{*}(X \setminus A)$ — внутренняя мера. Совокупность всех измеримых по Лебегу множеств обозначим через $\Sigma$.\\\textbf{Пример 7.1.} Покажем, что любое множество меры нуль измеримо по Лебегу и $\mu(A) = 0$. Действительно, если $A$ — множество меры нуль, то для любого $\varepsilon > 0$ существует система множеств $\{A_i\}_{i=1}^\infty$, $A_i \in K$, такая что $A \subset \bigcup_{i=1}^\infty A_i$ и $\sum_{i=1}^\infty m(A_i) < \varepsilon$. Тогда по определению внешней меры $\mu^{*}(A) \leq \sum_{i=1}^\infty m(A_i) < \varepsilon$, откуда $\mu^{*}(A) = 0$. Внутренняя мера $\mu_{*}(A) = m(X) - \mu^{*}(X \setminus A) \leq m(X)$, но также $\mu_{*}(A) \geq 0$. Однако из неравенства $\mu_{*}(A) \leq \mu^{*}(A)$ и равенства $\mu^{*}(A) = 0$ следует $\mu_{*}(A) = 0$. Таким образом, $\mu^{*}(A) = \mu_{*}(A) = 0$, и по упражнению 7.1 множество $A$ измеримо. Если множество $A$ измеримо, то его мерой полагают внешнюю меру, т.е. $\mu(A) = \mu^{*}(A)$.
\\\textbf{Теорема 7.1 (Критерий измеримости).} Пусть задано множество $X$ и алгебра $K$ с $\sigma$-аддитивной конечной мерой $m$. Тогда для любого множества $A \subset X$ следующие утверждения эквивалентны: 1) $A$ измеримо по Лебегу относительно меры $m$; 2) Для любого $\varepsilon > 0$ существует элементарное множество $B \in K$ такое, что $\mu^{*}(A \triangle B) < \varepsilon$. \textbf{Доказательство.} Докажем импликацию $1) \Rightarrow 2)$. Пусть $A$ измеримо, т.е. выполнено $\mu^{*}(A) + \mu^{*}(X \setminus A) = m(X)$. По определению внешней меры для любого $\varepsilon > 0$ существуют покрытия: — Для $A$: найдётся система множеств $\{B_i\}_{i=1}^\infty$, $B_i \in K$, такая что $A \subset \bigcup_{i=1}^\infty B_i$ и $\sum_{i=1}^\infty m(B_i) < \mu^{*}(A) + \varepsilon/3$. — Для $X \setminus A$: найдётся система множеств $\{C_i\}_{i=1}^\infty$, $C_i \in K$, такая что $X \setminus A \subset \bigcup_{i=1}^\infty C_i$ и $\sum_{i=1}^\infty m(C_i) < \mu^{*}(X \setminus A) + \varepsilon/3$. Выберем $N$ достаточно большим, чтобы $\sum_{i=N+1}^\infty m(B_i) < \varepsilon/3$. Положим $B = \bigcup_{i=1}^N B_i \in K$. Тогда $A \triangle B = (A \setminus B) \cup (B \setminus A)$. — Для $A \setminus B$: имеем $A \setminus B \subset \bigcup_{i=N+1}^\infty B_i$, откуда $\mu^{*}(A \setminus B) < \varepsilon/3$. — Для $B \setminus A$: имеем $B \setminus A = B \cap (X \setminus A) \subset \bigcup_{i=1}^\infty (B \cap C_i)$, откуда $\mu^{*}(B \setminus A) \leq \sum_{i=1}^\infty m(B \cap C_i)$. Используя разложение $C_i = (C_i \setminus B) \cup (B \cap C_i)$ и оценки для $\sum m(B_i)$ и $\sum m(C_i)$, получаем $\mu^{*}(B \setminus A) < 2\varepsilon/3$. Таким образом, $\mu^{*}(A \triangle B) \leq \mu^{*}(A \setminus B) + \mu^{*}(B \setminus A) < \varepsilon$. Импликация $1) \Rightarrow 2)$ доказана. Теперь докажем импликацию $2) \Rightarrow 1)$. Пусть для любого $\varepsilon > 0$ существует $B \in K$ такое, что $\mu^{*}(A \triangle B) < \varepsilon$. Так как $B \in K$, то $\mu^{*}(B) = m(B)$ и $\mu^{*}(X \setminus B) = m(X \setminus B)$, причём $m(B) + m(X \setminus B) = m(X)$. Используя свойство внешней меры $|\mu^{*}(A) - \mu^{*}(B)| \leq \mu^{*}(A \triangle B)$ и аналогичное для дополнений, получаем $|\mu^{*}(A) + \mu^{*}(X \setminus A) - m(X)| < 2\varepsilon$. В силу произвольности $\varepsilon$ заключаем, что $\mu^{*}(A) + \mu^{*}(X \setminus A) = m(X)$, т.е. множество $A$ измеримо.

%ВОРПОС: 11
\subsection*{11. Теорема о структуре множества измеримых по Лебегу множеств. Следствие об
измеримости счетного пересечения множеств.}
\textbf{Теорема 7.2.} Совокупность $\Sigma$ измеримых по Лебегу множеств образует $\sigma$-алгебру множеств, содержащую исходную алгебру $K$. \textbf{Доказательство.} Сначала докажем, что $\Sigma$ — алгебра. 1. Заметим, что в определении измеримости множества $A$ и $X \setminus A$ участвуют симметрично, поэтому если $A \in \Sigma$, то и $X \setminus A \in \Sigma$. 2. Пусть $A_1, A_2 \in \Sigma$. По теореме 7.1 для любого $\varepsilon > 0$ найдутся множества $B_1, B_2 \in K$ такие, что $\mu^{*}(A_1 \triangle B_1) < \frac{\varepsilon}{2}, \ \mu^{*}(A_2 \triangle B_2) < \frac{\varepsilon}{2}$. Положим $B = B_1 \cup B_2 \in K$. Тогда $(A_1 \cup A_2) \triangle B \subset (A_1 \triangle B_1) \cup (A_2 \triangle B_2)$, откуда $\mu^{*}((A_1 \cup A_2) \triangle B) \leq \mu^{*}(A_1 \triangle B_1) + \mu^{*}(A_2 \triangle B_2) < \varepsilon$. Следовательно, $A_1 \cup A_2 \in \Sigma$. 3. Так как $A_1 \cap A_2 = X \setminus ((X \setminus A_1) \cup (X \setminus A_2))$, то из уже доказанного следует, что пересечение двух измеримых множеств измеримо. 4. Разность $A_1 \setminus A_2 = A_1 \cap (X \setminus A_2)$ также измерима. Таким образом, $\Sigma$ — алгебра. Теперь докажем, что $\Sigma$ — $\sigma$-алгебра. Для этого достаточно показать, что счётное объединение измеримых множеств измеримо. Пусть $\{A_i\}_{i=1}^{\infty} \subset \Sigma$. Без ограничения общности можно считать, что множества $A_i$ попарно не пересекаются (иначе можно перейти к последовательности $B_1 = A_1$, $B_i = A_i \setminus \bigcup_{j=1}^{i-1} A_j$, которая состоит из попарно непересекающихся измеримых множеств, и $\bigcup_{i=1}^{\infty} A_i = \bigcup_{i=1}^{\infty} B_i$). Обозначим $A = \bigcup_{i=1}^{\infty} A_i$. Так как для любого конечного $n$ множество $\bigcup_{i=1}^{n} A_i$ измеримо (как конечное объединение измеримых), то по свойству монотонности внешней меры (которое верно и на $\Sigma$, так как внешняя мера совпадает с мерой на измеримых множествах) имеем: $\mu^{*}(A) \geq \mu^{*}\left(\bigcup_{i=1}^{n} A_i\right) = \sum_{i=1}^{n} \mu^{*}(A_i)$. Переходя к пределу при $n \to \infty$, получаем $\mu^{*}(A) \geq \sum_{i=1}^{\infty} \mu^{*}(A_i). $ Ряд справа сходится, так как его частичные суммы ограничены сверху числом $\mu^{*}(A) \leq m(X) < \infty$. Следовательно, для любого $\varepsilon > 0$ существует $N$ такое, что $\sum_{i=N+1}^{\infty} \mu^{*}(A_i) < \varepsilon$. Рассмотрим множество $C = \bigcup_{i=1}^{N} A_i \in \Sigma$. Тогда $\mu^{*}(A \triangle C) = \mu^{*}\left(\bigcup_{i=N+1}^{\infty} A_i\right) \leq \sum_{i=N+1}^{\infty} \mu^{*}(A_i) < \varepsilon$. По следствию 7.1 (или непосредственно по теореме 7.1, так как $C \in \Sigma$ и существует $B \in K$, близкое к $C$) множество $A$ измеримо. Таким образом, $\Sigma$ — $\sigma$-алгебра. Наконец, покажем, что $K \subset \Sigma$. Если $A \in K$, то по теореме 6.1 $\mu^{*}(A) = m(A)$ и $\mu^{*}(X \setminus A) = m(X \setminus A)$, откуда $\mu^{*}(A) + \mu^{*}(X \setminus A) = m(A) + m(X \setminus A) = m(X)$. Следовательно, $A \in \Sigma$. \\\textbf{Следствие 7.2.} Счётное пересечение измеримых множеств измеримо. \textbf{Доказательство.} Это следует из того, что $\bigcap_{i=1}^{\infty} A_i = X \setminus \bigcup_{i=1}^{\infty}(X \setminus A_i)$, а $\Sigma$ — $\sigma$


%ВОРПОС: 12
\subsection*{12. Теорема о свойстве сужения внешней меры на класс измеримых множеств. Полная
мера. sigma-конечная мера. Измеримость относительно sigma-конечной меры.}
\textbf{Теорема 7.3.} Сужение внешней меры $\mu^{*}$ на класс измеримых множеств $\Sigma$ задаёт счётно-аддитивную меру $\mu$, т.е. $\mu(A) = \mu^{*}(A)$ для всех $A \in \Sigma$. \textbf{Доказательство.} Сначала докажем конечную аддитивность. Пусть $A_1, A_2 \in \Sigma$ и $A_1 \cap A_2 = \varnothing$. По теореме 7.1 для любого $\varepsilon > 0$ найдутся множества $B_1, B_2 \in K$ такие, что $\mu^{*}(A_1 \triangle B_1) < \varepsilon, \ \mu^{*}(A_2 \triangle B_2) < \varepsilon$. Положим $B = B_1 \cup B_2 \in K$. Тогда $\mu^{*}(A \triangle B) \leq \mu^{*}(A_1 \triangle B_1) + \mu^{*}(A_2 \triangle B_2) < 2\varepsilon$, где $A = A_1 \cup A_2$. Используя свойство 6.5 внешней меры, получаем $|\mu^{*}(A) - \mu^{*}(B)| \leq \mu^{*}(A \triangle B) < 2\varepsilon$. Так как $B_1$ и $B_2$ могут пересекаться, вычислим меру $B$: $\mu(B) = m(B) = m(B_1) + m(B_2) - m(B_1 \cap B_2)$. Заметим, что $B_1 \cap B_2 \subset (A_1 \triangle B_1) \cup (A_2 \triangle B_2)$, поэтому $m(B_1 \cap B_2) \leq \mu^{*}(B_1 \cap B_2) \leq \mu^{*}(A_1 \triangle B_1) + \mu^{*}(A_2 \triangle B_2) < 2\varepsilon$. Кроме того, снова используя свойство 6.5, имеем $|\mu^{*}(A_1) - m(B_1)| \leq \mu^{*}(A_1 \triangle B_1) < \varepsilon, \ |\mu^{*}(A_2) - m(B_2)| < \varepsilon$. Подставляя эти оценки в формулу для $m(B)$, получаем $\mu^{*}(B) = m(B) > \mu^{*}(A_1) + \mu^{*}(A_2) - 3\varepsilon$. Тогда из предыдущей оценки следует $\mu^{*}(A) > \mu^{*}(A_1) + \mu^{*}(A_2) - 5\varepsilon$. В силу произвольности $\varepsilon$ получаем $\mu^{*}(A) \geq \mu^{*}(A_1) + \mu^{*}(A_2)$. С другой стороны, по свойству счётной полуаддитивности внешней меры имеем $\mu^{*}(A) \leq \mu^{*}(A_1) + \mu^{*}(A_2)$. Таким образом, $\mu^{*}(A) = \mu^{*}(A_1) + \mu^{*}(A_2)$. Итак, конечная аддитивность доказана. Теперь докажем счётную аддитивность. Пусть $\{A_i\}_{i=1}^{\infty}$ — последовательность попарно непересекающихся множеств из $\Sigma$, и $A = \bigcup_{i=1}^{\infty} A_i$. Для любого конечного $n$ из конечной аддитивности и монотонности следует: $\mu^{*}(A) \geq \mu^{*}\left(\bigcup_{i=1}^{n} A_i\right) = \sum_{i=1}^{n} \mu^{*}(A_i)$. Переходя к пределу при $n \to \infty$, получаем $\mu^{*}(A) \geq \sum_{i=1}^{\infty} \mu^{*}(A_i)$. С другой стороны, по свойству счётной полуаддитивности внешней меры, $\mu^{*}(A) \leq \sum_{i=1}^{\infty} \mu^{*}(A_i)$. Из этих двух неравенств следует равенство: $\mu^{*}(A) = \sum_{i=1}^{\infty} \mu^{*}(A_i)$. Таким образом, мера $\mu = \mu^{*}|_{\Sigma}$ счётно-аддитивна.
\\\textbf{Определение 7.2.} Мера $m$, заданная на алгебре $K$, называется \emph{полной}, если из условия $m(A)=0$ следует, что любое подмножество $B\subset A$ принадлежит $K$ и $m(B)=0$. \\\textbf{Определение 7.3.} Мера $\mu$, заданная на кольце $K\subset\mathcal P(X)$, называется \emph{sigma-конечной}, если существует последовательность множеств $\{A_n\}_{n=1}^{\infty}$, $A_n\in K$, такая, что $\mu(A_n)<\infty$ для всех $n$ и $X=\bigcup_{n=1}^{\infty}A_n$. \textbf{Определение 7.4.} Пусть мера $\mu$ на алгебре $K$ sigma-конечна. Множество $A\subset X$ называется \emph{измеримым}, если для каждого $n$ множество $A\cap X_n$ измеримо относительно сужения меры на $X_n$, где $\{X_n\}_{n=1}^{\infty}$ — разложение $X$ на попарно непересекающиеся множества конечной меры. При этом полагают $\mu(A)=\sum_{n=1}^{\infty}\mu(A\cap X_n). $

%ВОРПОС: 13
\subsection*{13. Измеримость по Лебегу, если $X$ – полуинтервал. Утверждения о мере Лебега
одноточечного множества, мере Лебега не более чем счетного ограниченного
множества, мере Лебега ограниченного промежутка, об измеримости по Лебегу
любого открытого или замкнутого множества и любого ограниченного борелевского
множества. Пример несчетного множества меры нуль.}
\textbf{Определение.} Множество $A$ называется измеримым по Лебегу на $[a,b)$, если выполнено равенство $\mu^{*}(A)+\mu^{*}([a,b)\setminus A)=b-a$. \\\textbf{Утверждение 8.1.} Множество, состоящее из одной точки, измеримо, и его мера равна нулю. \textbf{Доказательство.} Пусть $A=\{a\}$. Для любого $\varepsilon>0$ рассмотрим покрытие $[a,a+\varepsilon)$. Тогда $\mu^{*}(\{a\})\le\varepsilon$, откуда $\mu^{*}(\{a\})=0$. Следовательно, $\{a\}$ — множество меры нуль, а значит измеримо и $\mu(\{a\})=0$. \\\textbf{Утверждение 8.2.} Всякое не более чем счётное ограниченное множество точек прямой измеримо, и его мера равна нулю. \textbf{Доказательство.} Пусть $A=\{a_1,a_2,\dots\}$. Тогда $A=\bigcup_{i=1}^{\infty}\{a_i\}$. По утверждению 8.1 каждое $\{a_i\}$ измеримо и имеет меру нуль. Так как класс измеримых множеств является $\sigma$-алгеброй, то $A$ измеримо. По счётной аддитивности меры $\mu(A)=\sum_{i=1}^{\infty}\mu(\{a_i\})=0$. \\\textbf{Утверждение 8.3.} Любой ограниченный промежуток (интервал, полуинтервал, отрезок) измерим, и его мера равна его длине. \textbf{Доказательство.} Для полуинтервала $[\alpha,\beta)$ это следует из построения, так как он принадлежит исходному полукольцу $S$. Для отрезка $[\alpha,\beta]$ заметим, что $[\alpha,\beta]=[\alpha,\beta)\cup\{\beta\}$, и так как $\{\beta\}$ имеет меру нуль, то $\mu([\alpha,\beta])=\mu([\alpha,\beta))+\mu(\{\beta\})=(\beta-\alpha)+0=\beta-\alpha$. Для интервала $(\alpha,\beta)$ имеем $(\alpha,\beta)=[\alpha,\beta)\setminus\{\alpha\}$, откуда $\mu((\alpha,\beta))=\mu([\alpha,\beta))-\mu(\{\alpha\})=(\beta-\alpha)-0=\beta-\alpha$. Аналогично рассматриваются случаи полуинтервалов вида $(\alpha,\beta]$. \\\textbf{Утверждение 8.4.} Любое ограниченное открытое или замкнутое множество измеримо по Лебегу. \textbf{Доказательство.} Известно, что всякое ограниченное открытое множество $G\subset\mathbb{R}$ является объединением не более чем счётного числа попарно непересекающихся интервалов: $G=\bigcup_{i=1}^{\infty}(\alpha_i,\beta_i)$ (это следует, например, из того, что каждая точка открытого множества содержится в некотором максимальном интервале, и таких интервалов не более чем счётно). По утверждению 8.3 каждый интервал измерим, поэтому $G$ как счётное объединение измеримых множеств измеримо. Пусть теперь $F$ — ограниченное замкнутое множество. Выберем интервал $(\alpha,\beta)\supset F$. Тогда множество $G=(\alpha,\beta)\setminus F$ открыто (как разность открытого и замкнутого), следовательно измеримо. Тогда $F=(\alpha,\beta)\setminus G$ измеримо как разность двух измеримых множеств.

%ВОРПОС: 14
\subsection*{14. Пример неизмеримого по Лебегу множества. Измеримость по Лебегу неограниченных
множеств. Замечание об измеримости произвольного множества числовой прямой.
Пример неограниченного измеримого по Лебегу множества. Теорема о sigma-конечности
меры, определенной на sigma-алгебре всех измеримых по Лебегу множеств (без док-ва).}
\textbf{Определение 8.1 (Измеримость неограниченного множества).} Множество \(A\subset\mathbb{R}\) называется измеримым по Лебегу, если для любого натурального \(n\in\mathbb{N}\) ограниченное множество \(A\cap[-n,n)\) измеримо по Лебегу (в смысле меры на ограниченном промежутке). Мерой Лебега множества \(A\) называется предел (возможно, бесконечный) \(\mu(A)=\lim_{n\to\infty}\mu\bigl(A\cap[-n,n)\bigr)\). \textbf{Замечание 8.1.} Альтернативный подход: представим \(\mathbb{R}\) в виде объединения попарно непересекающихся полуинтервалов единичной длины \(\mathbb{R}=\bigcup_{n=-\infty}^{\infty}[n,n+1)\). Тогда множество \(A\subset\mathbb{R}\) измеримо тогда и только тогда, когда для каждого \(n\) измеримо множество \(A_n=A\cap[n,n+1)\) (относительно меры Лебега на этом отрезке). При этом \(\mu(A)=\sum_{n=-\infty}^{\infty}\mu(A_n)\), где, если ряд расходится, то полагают \(\mu(A)=\infty\). \textbf{Пример 8.3.} Рассмотрим неограниченное множество \(A=\bigcup_{n=1}^{\infty}\bigl(n,\;n+\tfrac{1}{3^{n}}\bigr)\). Оно измеримо, так как каждое \(A_n=A\cap[n,n+1)=(n,n+\tfrac{1}{3^{n}})\) измеримо. Его мера равна \(\mu(A)=\sum_{n=1}^{\infty}\mu(A_n)=\sum_{n=1}^{\infty}\tfrac{1}{3^{n}}=\tfrac{1}{2}\). \textbf{Теорема 8.1.} Совокупность \(\Sigma\) всех измеримых по Лебегу подмножеств \(\mathbb{R}\) является \(\sigma\)-алгеброй; \(\mu\) является \(\sigma\)-конечной мерой на этой \(\sigma\)-алгебре.

%ВОРПОС: 15
\subsection*{15. Скачок функции. Утверждение об оценке суммы скачков неубывающей функции на
отрезке. Утверждение о множестве точек разрыва монотонно неубывающей на
отрезке функции}
\textbf{Пусть} $F:[a,b]\to\mathbb{R}$ — монотонно неубывающая функция. Тогда в любой точке $x_0\in(a,b)$ существуют односторонние пределы: $F(x_0-0)=\lim_{x\to x_0-}F(x)=\sup\{F(x):a\leq x<x_0\}$, $F(x_0+0)=\lim_{x\to x_0+}F(x)=\inf\{F(x):x_0<x\leq b\}$. В точках $a$ и $b$ существуют соответственно $F(a+0)$ и $F(b-0)$. Очевидно, что $F(x_0-0)\leq F(x_0)\leq F(x_0+0)$. Если $F(x_0-0)=F(x_0+0)$, то функция непрерывна в точке $x_0$. В противном случае точка $x_0$ называется \emph{точкой разрыва первого рода}. Разность $\Delta F(x_0)=F(x_0+0)-F(x_0-0)$ называется \emph{скачком функции} в точке $x_0$. Для концов отрезка полагают $\Delta F(a)=F(a+0)-F(a)$, $\Delta F(b)=F(b)-F(b-0)$. Если $F(x_0)=F(x_0-0)$, то функция называется \emph{непрерывной слева} в точке $x_0$; если $F(x_0)=F(x_0+0)$, то \emph{непрерывной справа}. \\\textbf{Утверждение 9.1.} Пусть $F$ — неубывающая функция на отрезке $[a,b]$ и $c_1,\dots,c_n$ — любые точки этого отрезка. Тогда $\sum_{i=1}^{n}\Delta F(c_i)\leq F(b)-F(a)$. \textbf{Доказательство.} Без ограничения общности можно считать, что точки упорядочены: $a\leq c_1<c_2<\dots<c_n\leq b$. Рассмотрим сумму скачков в этих точках. Заметим, что $\sum_{i=1}^{n}\Delta F(c_i)=(F(c_1+0)-F(c_1-0))+\dots+(F(c_n+0)-F(c_n-0))$. Для внутренних точек $c_i$ ($1<i<n$) имеем $F(c_i-0)\leq F(c_{i-1}+0)$ и $F(c_i+0)\leq F(c_{i+1}-0)$ в силу монотонности. Поэтому, группируя слагаемые, получаем $\sum_{i=1}^{n}\Delta F(c_i)\leq (F(c_1+0)-F(a))+(F(c_2+0)-F(c_1-0))+\dots+(F(b)-F(c_n-0))\leq F(b)-F(a)$, поскольку все разности в скобках неотрицательны и не превосходят соответствующих приращений функции. \\\textbf{Утверждение 9.2.} Множество точек разрыва монотонно неубывающей на отрезке $[a,b]$ функции не более чем счётно. \textbf{Доказательство.} Для каждого натурального $k$ рассмотрим множество $M_k=\{x\in[a,b]:\Delta F(x)\geq \tfrac{1}{k}\}$. Из утверждения 9.1 следует, что в множестве $M_k$ не более $k(F(b)-F(a))$ точек, т.е. $M_k$ конечно. Поскольку каждая точка разрыва принадлежит некоторому $M_k$ (скачок положителен), то множество всех точек разрыва есть объединение счётного семейства конечных множеств $\bigcup_{k=1}^{\infty}M_k$, а значит, не более чем счётно.

%ВОРПОС: 16
\subsection*{16. Производные числа Дини. Пример. Утверждение о производной монотонно
неубывающей на отрезке функции (без док-ва). Функция скачков.}
Введём понятие производных чисел Дини. Для этого рассмотрим предел отношения \(\frac{f(x)-f(x_0)}{x-x_0}\) при \(x \to x_0\). Данный предел может не существовать, но всегда существуют конечные или бесконечные верхние и нижние пределы: \(\Lambda_{\text{пр}} = \limsup_{x \to x_0+0} \frac{f(x)-f(x_0)}{x-x_0}\). Величина \(\Lambda_{\text{пр}}\), определяемая этим соотношением, называется правым верхним производным числом. Аналогично, \(\lambda_{\text{пр}} = \lim_{x \to x_0+0} \frac{f(x)-f(x_0)}{x-x_0}\), \(\Lambda_{\text{лев}} = \lim_{x \to x_0-0} \frac{f(x)-f(x_0)}{x-x_0}\), \(\lambda_{\text{лев}} = \lim_{x \to x_0-0} \frac{f(x)-f(x_0)}{x-x_0}\). Здесь \(\lambda_{\text{пр}}\) — нижнее правое производное число, \(\Lambda_{\text{лев}}\) — верхнее левое, а \(\lambda_{\text{лев}}\) — нижнее левое производные числа. Ясно, что всегда \(\lambda_{\text{пр}} \leq \Lambda_{\text{пр}}, \ \lambda_{\text{лев}} \leq \Lambda_{\text{лев}}\). \\\textbf{Утверждение 9.3.} Монотонно неубывающая на отрезке \([a,b]\) функция имеет почти всюду (относительно меры Лебега) конечную производную.\\Важным классом монотонных функций являются функции скачков. Пусть задано конечное или счётное множество точек \(\{x_n\} \subset [a,b]\) и соответствующие положительные числа \(h_n\) такие, что ряд \(\sum h_n\) сходится. Определим функцию \(F(x) = \sum_{x_n < x} h_n\), где сумма берётся по всем индексам \(n\), для которых \(x_n < x\). Дополнительно полагаем \(F(a)=0\) (если ни одна точка \(x_n\) не меньше \(a\), то сумма пустая, т.е. равна 0). Функция \(F\) монотонно неубывающая, ограниченная (поскольку ряд сходится) и непрерывная слева в каждой точке. Точками разрыва функции \(F\) являются именно точки \(x_n\), причём скачок в точке \(x_n\) равен \(h_n\): \(\Delta F(x_n) = F(x_n+0) - F(x_n-0) = h_n\). Такая функция называется функцией скачков.





\end{document}
